\subsection{箭图的连通分支}

设 \(C = (S, F, o, t)\) 为箭图。考虑 S 上的等价关系 \(R_{C}\) ，即使得对 C 的每个箭 f 关系 \(\mathrm{R}_{\mathrm{C}}\{o(f), t(f)\}\) 成立的最细者。C 的两个顶点 a、b 对这个关系等价，当且仅当存在整数 \(n \geqslant 0\) 、C 的顶点 \(a_{0}, \ldots, a_{n}\) 与 C 的箭 \(f_{1}, \ldots, f_{n}\) ，使得 \(a_{0} = a\) ，\(a_{n} = b\) ，且对 \(1 \leqslant i \leqslant n\) ，箭 \(f_{i}\) 或者连接 \(a_{i-1}\) 与 \(a_{i}\) ，或者连接 \(a_{i}\) 与 \(a_{i-1}\) 。

这个等价关系的等价类称为 C 的连通分支。以 \(\pi_0(\mathrm{C})\) 表示 C 的连通分支所成的集。最后，若 C 至多有一个连通分支，则称 C 是连通的。

设 \(\varphi\colon\mathrm{C}\to\mathrm{C}^{\prime}\) 为箭图态射。映射 \(\operatorname{Som}(\varphi)\) 通过过渡到商，定义了一个从 \(\pi_{0}(\mathrm{C})\) 到 \(\pi_{0}(\mathrm{C}^{\prime})\) 的映射，我们把它记作 \(\pi_{0}(\varphi)\) 。

